\clearpage
\properlane{grace-and-willing-service}{The received gift becomes willing service}
The second reading begins with a difficulty inside the worshipper's action. If God gives strength, grace and even the capacity to will, what remains for a person to do? Conversely, if faith works and love labours, why should thanksgiving credit God? Cyrus's unknowing reception of power, Paul's grateful recognition of a fruitful community, and the Collect's petition concerning the will make these questions concrete. Aquinas explains how divine action reaches willing itself; Chrysostom explains why the resulting human life must be active, costly and sincere.

\subsection{A gift can be used without its giver being known}
Cyrus is not merely an example of someone else's politics. His condition reveals a possibility within every gifted life: capacity can exceed acknowledgment. Aquinas distinguishes the benefit received by the people from the recognition sought in Cyrus. A good can really occur through a person who does not yet know its source. That distinction prevents a community from treating its own usefulness as sufficient evidence of conversion. It may accomplish something valuable while misunderstanding whom its accomplishment serves.\footnote{Aquinas, \work{Super Isaiam}, chapter 45, paragraph 85005.}

Jerome deliberately brings this problem into ecclesial life. At the end of his exposition he applies it to the teacher given speech and wisdom. Such a teacher must refer the glory to the giver rather than imagine that his power of speaking originates in himself. Otherwise the rebuke addressed to Cyrus can become his own. Jerome's move is searching because the teacher's task is religious. Knowing how to speak about God does not automatically mean acknowledging God through the manner of one's life.\footnote{Jerome, \work{Commentary on Isaiah}, XII, Isaiah 45:1--7, pp.128--129, concluding tropological application.}

The Entrance and Collect answer this condition through petition. The Entrance's confident cry grows, in Augustine's explanation, from prior divine hearing. It begins with help rather than achievement. The Collect asks concerning the will's devotion, placing the very readiness to serve within dependence upon God. These prayers do not deny the assembly's responsibility. They name its deepest need: not only tools for good work, but a will capable of doing that work sincerely.

\subsection{God acts within willing}
The acclamation's two short phrases stand within Philippians 2, where exhortation to work out salvation is immediately joined to God's action in willing and doing. Aquinas treats that wider context in lectio 3. He excludes the notion that free choice can save without divine help. He also excludes the denial of free choice through fate or providence: God's action takes place inwardly in the will, moving the person to act. Divine help is not a second human agent pushing the person aside.\footnote{Aquinas, \work{Super Philippenses}, 2, lectio 3, paragraph 87838, exposition of 2:12--13; these verses are the acclamation's context, not its appointed text.}

His next distinction reaches the beginning of action. It is insufficient to assign willing to ourselves and leave God only the completion of what we independently decide. The text says both willing and doing. Nor is this help payment for merit earned before grace. Aquinas's account therefore does more than add divine assistance to a programme whose source remains wholly our own. The good decision itself can be truly ours and truly the fruit of grace.

This is why petition for the will is intelligible. If willing were wholly independent, prayer could ask only for favourable circumstances. If willing were unreal, exhortation and response would have no subject. Aquinas preserves both dependence and agency. The consequence is neither self-congratulation nor passivity, but a readiness to receive help at the point where a person chooses. The Collect's concern is not completed when a person has identified a duty; it reaches the manner in which he becomes willing to perform it.

Lectio 4 brings this inward account to the acclamation's outward image. The saints shine by participation in the Word's light. Receiving does not extinguish their activity any more than receiving light makes a visible body disappear. Their distinct lives become the place where what was received is seen. Holding Christ's life-giving word supplies the source and direction of that visibility. This keeps the acclamation from being reduced to the advice to become impressive.

\subsection{Thanksgiving recognizes action instead of replacing it}
Paul's opening thanksgiving provides the clearest human instance. He credits God for a people whose faith has work, whose love bears labour and whose hope sustains endurance. Chrysostom explains that their conduct gives grounds for thanksgiving; thanksgiving in turn teaches humility. The two acts are not competitors. To acknowledge that a gift made an action possible is not to say the action did not occur. To recognize the action is not to say the gift was unnecessary.\footnote{Chrysostom, \work{1 Thessalonians}, I, on 1:2--5.}

Chrysostom makes the words costly. Steadfastness has to survive what might turn it aside. Love is not exhausted by an affectionate name; it endures labour. Hope has to remain during affliction rather than wait for affliction to end before calling itself hope. This explains why the community's good cannot be measured merely by ease or prosperity. The gift is visible in sustained fidelity where favourable conditions are lacking.

Aquinas's account of election adds a necessary order. Fruits manifest God's choice; they do not retroactively purchase it. His alternative explanations of the Gospel's power and fullness allow attention to the preacher's conduct and the completeness of instruction as well as to signs. An assembly seeking evidence of grace should therefore attend to the ordinary but demanding fruits Paul names: truthfulness under pressure, charity that continues when inconvenient, and hope that survives delay. These are applications of the text's criteria, not a method for declaring the hidden destiny of particular people.\footnote{Aquinas, \work{Super I Thessalonicenses}, 1, lectio 1, paragraph 87862.}

\subsection{From truthful words to a sincere offering}
The Gospel's questioners supply a contrary example. They speak truth about Jesus while trying to entrap him through it. Correct words can be detached from a truthful intention. That disjunction gives the Collect's sincere service a practical edge: the mouth can render honour while the will attempts to secure control. Jesus' exposure of the trap calls the speaker into the truth he has spoken.

Aquinas's mystical return of God's image asks more than an outward transfer. The person must answer the giver. Yet the literal tribute question remains: a claim to inward sincerity cannot excuse injustice in practical obligations. This is where the first reading of the formulary corrects a danger in the second. Grace does not make the tax disappear into an allegory of the soul. The person transformed inwardly must act truthfully toward neighbours and within the civil life in which they live.

The responsorial Psalm gives that transformed response a voice. Augustine's preacher declares God's glory instead of his own. The shift is small in grammar and large in life: a gift becomes praise when the giver receives the acknowledgment. At the altar, the Prayer over the Offerings adds that sincere service still needs purification. A willingness once awakened is not a permanent exemption from mixed motives. Serving the mysteries becomes an occasion for the servant's renewal rather than a reward for having already become complete.

\subsection{Mercy nourishes; the Son first gives himself}
The Psalm Communion option gathers fear and hope around the Lord's mercy. Aquinas's explanation is exact: fear alone can lead toward despair; hope severed from fear can become presumption. Together they describe a person who takes both sin and mercy seriously. Bodily food sustains bodily life; the Word sustains spiritual life; sacramental nourishment belongs within the same account of God's strengthening help. The recipient does not need to pretend to be full before being fed.\footnote{Aquinas, \work{Super Psalmos}, 32, §16; Augustine, \work{Expositions}, XXXII, §§18--22.}

Augustine's patient soul is therefore a fitting companion to Paul's enduring community. Need can persist while mercy is real. Neither author gives a warrant for blaming every sufferer for deficient faith. The hope in the Psalm reaches incorruptible nourishment; it cannot be tested simply by asking whether present circumstances have become comfortable. This makes perseverance a response to a promised good, rather than a denial of actual hunger.

The Mark option sharpens the order of gift, reception and imitation. Bede's explanation begins with the Son taking the servant's form and shedding his blood for the world. His gloss of ``for many'' then names those willing to believe. The recipient's act has a specified object: the saving gift of Christ. Believing is not self-redemption, and Christian service does not supply a missing ransom equal to Christ's own. Before the disciple gives himself, he receives the one who has given himself.\footnote{Bede, \work{In Marci}, III, Mark 10:45, PL 92, col.236B.}

Bede's reference to willing belief does not assign an independent origin to faith's will. Aquinas's account of God's inward action supplies that distinct question's answer. Chrysostom's exposition of Matthew's parallel then develops the consequence for conduct: Christ's abasement unmakes the disciples' ambition. A person can imitate the pattern of service without becoming another redeemer. The priority of Christ's act is what gives imitation its Christian shape, rather than leaving it a contest in moral heroism.

The Prayer after Communion keeps this life open. Present participation prepares for eternal gifts; it does not convert progress into possession of the end. The assembly can leave with work to do because grace has not made work pointless. It can leave without claiming mastery because its work is not the source of salvation. Readiness, purification and growth belong to a single received life whose fulfilment still calls it onward.

\subsection{The four senses}
\begin{fourSenses}
\item[Literal.] Cyrus receives strength without corresponding knowledge; Paul recognizes an active congregation and thanks God for it; Philippians joins received life to visible conduct. The prayers ask for the will's disposition, the servant's purification and sacramental progress. The Gospel exposes truthful language used with an untruthful purpose.
\item[Allegorical.] Christ's self-gift gives life to a body whose light comes from him. Augustine's Christ-and-Church prayer, Aquinas's participated light and sacramental nourishment, and Bede's believing reception locate the community's action in a prior gift. Cyrus remains a limited type, not a Christ who fails to know the Father.
\item[Moral.] Acknowledge the giver within one's abilities, choose and carry out good work with received help, and let service expose what still needs purification. Hold fear together with hope. The virtues Paul names become strenuous love and sustained fidelity, without making effort a purchase of election.
\item[Anagogical.] Hope acts toward a good not yet fully possessed. Present nourishment sustains patient life on the way to eternal gifts; progress does not end the pilgrimage. Final fulfilment gives present work a direction beyond its immediate success or reward.
\end{fourSenses}
